% NOETHER PAPER 6 — KOREAN TRANSLATION-PRODUCER DRAFT — UNCHECKED
% Unit: T05-U43; current-authority whole lines 4954--4970
% Source slice: 1,113 LF UTF-8 bytes; SHA-256 A5C58E98E7EDA9CD0FD85ECEC09DE74A58ADC6BB157F3347A293C482A073365C
% Pointer: NOETH-DE-AUTH-v007-20260804; SHA-256 A6A8FC8E5AC24ACAF49DFD55B4B58FA3DA882EF8C3FDD4D136220C8751045156
% German authority: NOETH-DE-ED-0001; 2,153,565 raw bytes; SHA-256 D1F06B311F6CBD991DD247D745DD9A72DDE326A20396DF43CFE0C8EDB1593CDB
% Producer translation only: no source/Korean/formula review, compilation, or rendering.
% Sense window: algebraischer Körper über is rendered as 대수적 확대체, while Körper alone remains 체(體).
% Terminology debt: Adjunktion→첨가/부가 and Zwischenkörper→중간체/중간 체 remain checker-held.
% Standardization hold: ko-KR 첨가·중간체 and Hanja 체(體) are provisional; ko-KP equivalents remain unresolved.

\[
x_1,x_2,\ldots,x_\rho
\]
가
\[
\Omega\bigl(g_1^*(x)\cdots g_\rho^*(x)\bigr)
\]
에 관하여 대수적 원소라는 것이다. 따라서 $x_1\cdots x_\rho$를 첨가하여 생기는 체(體)
\[
\Omega\bigl(g_1^*\cdots g_\rho^*,x_1\cdots x_\rho\bigr)=\Omega(x_1\cdots x_\rho)
\]
는 $\Omega(g_1^*\cdots g_\rho^*)$ 위의 (유한) 대수적 확대체이다. 알려진 정리\srcfn{**)}{Vgl. etwa Weber, Lehrbuch d. Algebra 1, \S\ 144 (1. Aufl.) oder \S\ 55 (kleine Ausgabe).}에 따라 $\Kfield^*_{\rho\rho}$는 $\Omega(g_1^*\cdots g_\rho^*)$와 $\Omega(x_1\cdots x_\rho)$의 중간체로서 그 자체가 $\Omega(g_1^*\cdots g_\rho^*)$ 위의 대수적 확대체이며, 한편 $\Omega(x_1\cdots x_\rho)$는 $\Kfield^*_{\rho\rho}$ 위의 대수적 확대체이다. 그러므로 $\Kfield^*_{\rho\rho}$는 $\Omega(g_1^*\cdots g_\rho^*)$에 $\Kfield^*_{\rho\rho}$의 유한 개 원소를 첨가하여 얻어지며, 또는 적절히 택한 함수 \emph{하나}를 첨가하여도 얻어진다. 즉,
\[
\Kfield^*_{\rho\rho}
=\Omega(g_1^*\cdots g_\rho^*,g_{\rho+1}^*\cdots g_k^*)
=\Omega(g_1^*\cdots g_\rho^*g_0^*).
\]
